% =====================================================================
\section{The system and the joint controller}
\label{sec:system}
% =====================================================================

\subsection{What was built}
\label{sec:built}

\polyforge{} comprises a Go control plane and gateway with a per-tenant
semantic cache, a Python planner service, and a Kubernetes operator that
applies the planner's decisions under guardrails (move clamps, a plan
deadline, and a hold-last-good-plan fallback). The gateway's replicas are
the tenants' serving pods: they authenticate, cache, meter and forward each
request; the models run outside the cluster as endpoints with a price and a
latency, so no replica is a GPU model server. The planner service runs the
same controller code as the simulator in which every campaign of this
paper ran (Section~\ref{sec:method}); the operator path was exercised on a
local Kubernetes-in-Docker cluster and on one rented GPU host
(Section~\ref{sec:live}), but none of the comparisons reported in
Sections~\ref{sec:chain}--\ref{sec:width} ran on a cluster. Several
components in the repository (an analytics sink, OIDC and secret-store
integrations, a service-mesh configuration and a network policy) were
never part of a measured run, and we make no claim about them.

\subsection{The planner}
\label{sec:jcac}

Every 10\,s, \jcac{} chooses a configuration $x_i=(r_i,c_i,\tau_i)$ for
each tenant $i$: a replica count, a cache budget and a model tier. The
choice is restricted to a \emph{lattice} of moves anchored at the
interval-start state: $\Delta r\in\{-2,\dots,2\}$, a move of at most one
cache level among $\{0,64,128,256,512,1024\}$\,MB, and
$\tau\in\{\text{none},\text{small},\text{mid},\text{large}\}$, where
``none'' makes no model call. That gives at most 60 candidates per tenant.
Over a six-step (60\,s) horizon the planner minimises
\begin{equation}
\label{eq:plan}
\Phi(x)=\frac{\alpha}{c_0}\sum_i \mathrm{cost}_i
 +\beta\sum_i\log\bigl(1+\mathrm{exc}_i\bigr)
 +\gamma\bigl(1-\mathrm{Jain}(s)\bigr)
 +\lambda\sum_i \mathrm{sw}_i ,
\end{equation}
where $\mathrm{cost}_i$ is the projected spend of tenant $i$ (replica time,
cache memory and the tier's price per cache miss), $\mathrm{exc}_i$ is the
traffic-weighted relative overshoot of its p95 latency over its service
level objectives (SLOs), $s_i=1-v_i$ with $v_i\in[0,1]$ the violation (the
overshoot capped at 1 per request family) enters Jain's fairness
index~\cite{jain1984}, and $\mathrm{sw}_i$ counts changed knobs. The weights
are $\alpha=1$, $\beta=2$, $\gamma=0.5$, $\lambda=0.05$, and
$c_0=\$0.01$ per tenant-step. An interference penalty, zero unless a tenant
is flagged as a noisy neighbour, is omitted because no campaign here flags
one.

Two hard filters precede the objective. A candidate whose projected spend
exceeds the tenant's hourly dollar budget, pro-rated to the interval, is
discarded, and the tenants' combined replicas and cache may not exceed the
cluster's caps. If no candidate survives, the tenant falls back to its
minimum replicas, no cache and tier ``none''. We call a tenant
\emph{shed} on a step whose tier is ``none'', whether the objective chose it
or the budget forced it. Shedding is the controller's designed response to
an unaffordable interval; no comparator in this study can shed, and the
objective prices a shed only through its SLO terms
(Section~\ref{sec:objective}).

The planner holds each candidate fixed over the horizon, on a forecast that
extends each tenant's linear trend over its last three intervals, applies
only the first move, and re-plans every interval: receding-horizon control
with move blocking. The solver runs block-coordinate descent over tenants.
Each tenant best-responds exactly over its own lattice, within the cluster
caps that tenants earlier in the sweep leave free; sweeps accept only
strict improvements of $\Phi$ and repeat until no tenant changes, which
terminates because the lattice is finite. With uniform fairness weights,
$\Phi$ is an exact potential~\cite{monderer1996} of the game in which each
tenant chooses its own configuration, so the sweeps end at a
coordinate-wise equilibrium: no tenant can lower $\Phi$ by changing only its
own configuration. This is joint search \emph{within} a tenant, not a
global optimum across tenants, and the result can depend on the sweep
order. An earlier, two-sweep version of the solver could move a tenant past
its clamps on the second sweep (up to $\pm4$ replicas); every campaign
reported here uses the anchored, convergent solver, which never reached its
50-sweep safety cap where tracked. Planning 8 tenants takes a median of
31\,ms (p95 44\,ms) per cycle on one laptop core.

\subsection{The ablations that isolate coordination}
\label{sec:ablations}

Two ablations share everything with \jcac{}---objective, forecast, budget
filter, cluster caps and lattice---and change only the search over one
tenant's knobs. The \emph{decoupled} arm best-responds each knob with the
other two held at the incumbent and combines the three winners; if the
combined move breaks a budget or a cap, it keeps the best feasible subset of
the individual moves. The \emph{knob-wise} arm takes one-knob best responses in turn
until no knob changes. The difference between \jcac{} and
the decoupled arm is therefore the value of searching the knobs jointly,
with nothing else changed.
